\scanpage{8}
\mspage{7}
\[
F\mid M^{+0}(S)\simeq F_\xi
\]
où $F_\xi$ est le foncteur évident sur $M^{+0}(S)$ : fibre en $\xi$. (NB. $S$ étant connexe, ces foncteurs sont isom entre eux, donc la \uline{condition} d'association ne dépend pas de $\xi$. Mais il faut \unsure{le|la} considérer plutôt comme une \uline{donnée} supplémentaire, comme \uline{donnée} de l'isom en question).

\uline{Ex 1} Si $\operatorname{source}\xi=\operatorname{Spec}(\mathbb C)$, \deleted{alors} $F=F^{\mathrm{Betti}}_\xi$, alors on a une telle donnée. \inserted{\uline{Ex. 2} $T^\ell_\xi$}

\uline{Ex} \deleted{\unsure{2}}3 Si $S=\operatorname{Spec}k$, $k$ corps de car. nulle (p. ex. $k=Q$) et si $F$ est le foncteur fibre \deleted{sur}\inserted{à coefficients dans} $k$ $T_{\mathrm{DR}}$ \deleted{alors}, ou $T_{\mathrm{Hdg}}$, alors il n'est pas en général associé aux pts géométriques.

\uline{Ex} 4 Si \deleted{$k$ est} $S$ est le spectre d'un corps alg clos, alors tt $F$ est associé à $\xi$, du moins si $k$ est dénombrable \inserted{i.e. de degré de transc. absolu dénombrable}, donc $M(k)$ à engendrement dénombrable\ldots .

\deleted{Alors on trouve un isom canonique
\[
\pi_1(S,\xi)\simeq\varprojlim G_\alpha/G_\alpha^0
\]
\inserted{Les $G_\alpha/G_\alpha^0$ sont des groupes constants sur $K$ \ldots. :}
Cela résulte du \deleted{fait que} $\otimes$-foncteur canonique
\[
\mathrm{Mod}\unsure{\text{ adm}|\text{ disc}|\text{ dis}|\text{ cont}}(\pi_1)\simeq M^{+0}(S)\longrightarrow M(S)
\]
qui définit
\[
G\longrightarrow\pi_1
\]}
% archive sheet number: 7

